% =====================================================================
%  Thème 2 — EXERCICES — Niveau FACILE
% =====================================================================
\documentclass[11pt,a4paper]{article}
\usepackage[utf8]{inputenc}
\usepackage[T1]{fontenc}
\usepackage{lmodern}
\usepackage[french]{babel}
\usepackage{amsmath,amssymb,mathtools}
\usepackage{icomma}
\usepackage{xcolor}
\usepackage[most]{tcolorbox}
\usepackage{enumitem}
\usepackage{array}
\usepackage{booktabs}
\usepackage{fancyhdr}
\usepackage[hidelinks]{hyperref}
\usepackage[a4paper,margin=2.2cm,headheight=26pt,headsep=10pt]{geometry}

\definecolor{primary}{RGB}{223,87,99}
\definecolor{primarydark}{RGB}{150,42,55}

\tcbset{boxbase/.style={enhanced,breakable,boxrule=0.9pt,arc=2.5pt,
   left=3mm,right=3mm,top=2.3mm,bottom=2.3mm,
   fonttitle=\bfseries,coltitle=white,toptitle=0.6mm,bottomtitle=0.6mm}}
\newtcolorbox[auto counter]{exercice}[1][]{boxbase,colback=primary!4,colframe=primary,
  title={\textsc{Exercice}~\thetcbcounter\ifx\relax#1\relax\relax\else\textmd{ --- \itshape #1}\fi}}

\pagestyle{fancy}
\fancyhf{}
\fancyhead[L]{\small\textcolor{primarydark}{\textbf{Fiche 6} — Les limites}}
\fancyhead[R]{\small\textcolor{primarydark}{\textsc{Thème 2 — Facile}}}
\fancyfoot[C]{\small\thepage}
\renewcommand{\headrulewidth}{0.8pt}
\renewcommand{\headrule}{\hbox to\headwidth{\color{primary}\leaders\hrule height \headrulewidth\hfill}}

\newcommand{\R}{\mathbb{R}}

\begin{document}

\begin{center}
{\Large\bfseries\color{primarydark}Thème 2 — Règle $\dfrac{a}{0^{\pm}}$ \& asymptote verticale}\\[2pt]
{\large\color{primary}Exercices — Niveau Facile}
\end{center}
\vspace{4pt}
{\small\itshape But : appliquer la règle des signes $\frac{a}{0^{\pm}}\to\pm\infty$ et lire le signe
d'un dénominateur autour d'un point.}
\vspace{6pt}

\begin{exercice}[appliquer la règle des signes]
Donner le résultat ($+\infty$ ou $-\infty$) de chaque forme.
\begin{enumerate}[label=\alph*),leftmargin=2em,itemsep=3pt]
  \item $\dfrac{5}{0^{+}}$ \qquad b) $\dfrac{-3}{0^{+}}$ \qquad c) $\dfrac{2}{0^{-}}$ \qquad d) $\dfrac{-7}{0^{-}}$
\end{enumerate}
\end{exercice}

\begin{exercice}[signe du dénominateur]
Pour chaque expression, dire si le dénominateur tend vers $0^{+}$ ou $0^{-}$ du côté indiqué.
\begin{enumerate}[label=\alph*),leftmargin=2em,itemsep=3pt]
  \item $x-4$ quand $x\to 4^{+}$ \qquad b) $x-4$ quand $x\to 4^{-}$
  \item $2-x$ quand $x\to 2^{+}$ \qquad d) $(x-1)^2$ quand $x\to 1^{-}$
\end{enumerate}
\end{exercice}

\begin{exercice}[limites latérales élémentaires]
Calculer les deux limites latérales.
\begin{enumerate}[label=\alph*),leftmargin=2em,itemsep=3pt]
  \item $\displaystyle\lim_{x\to 0^{+}}\frac{3}{x}$ \qquad et \qquad $\displaystyle\lim_{x\to 0^{-}}\frac{3}{x}$
  \item $\displaystyle\lim_{x\to 4^{+}}\frac{1}{x-4}$ \qquad et \qquad $\displaystyle\lim_{x\to 4^{-}}\frac{1}{x-4}$
\end{enumerate}
\end{exercice}

\begin{exercice}[une fonction rationnelle simple]
Soit $f(x)=\dfrac{7}{x+3}$, de domaine $\R\setminus\{-3\}$.
\begin{enumerate}[label=\alph*),leftmargin=2em,itemsep=3pt]
  \item Le numérateur tend-il vers $0$ en $-3$ ? Quel est son signe ?
  \item Calculer $\displaystyle\lim_{x\to -3^{+}} f(x)$ et $\displaystyle\lim_{x\to -3^{-}} f(x)$.
\end{enumerate}
\end{exercice}

\begin{exercice}[repérer une asymptote verticale]
Pour chaque fonction, donner la valeur $x_0$ qui annule le dénominateur (sans l'annuler au
numérateur) : c'est l'abscisse d'une asymptote verticale.
\begin{enumerate}[label=\alph*),leftmargin=2em,itemsep=3pt]
  \item $\dfrac{x+1}{x-5}$ \qquad b) $\dfrac{4}{2x+6}$ \qquad c) $\dfrac{5}{x^2}$
\end{enumerate}
\end{exercice}

\end{document}
